\documentclass[11pt,letterpaper]{article}
\usepackage[margin=0.7in,headheight=15pt]{geometry}
\usepackage[T1]{fontenc}\usepackage{lmodern,amsmath,amssymb,fancyhdr,parskip,hyperref}
\hypersetup{colorlinks=true,urlcolor=blue}
\pdfinfoomitdate=1\pdftrailerid{}\pdfsuppressptexinfo=15
\pagestyle{fancy}\fancyhf{}\lhead{Week 80 / Midnight laboratory}\rhead{Adult guide}\cfoot{Bellingham Math Circle / Week 80 / M80-FAC-v2 / \thepage}
\setlength{\parskip}{7pt}\setlength{\parindent}{0pt}
\begin{document}
\section*{Mathematical overview}
The light starts OFF at 11:00 p.m., switches ON at 11:30, OFF at 11:45, ON at 11:52:30, and so on. Normalize 11:00 p.m. to time 0 and midnight to time 1. The switch times $t_n=1-2^{-n}$ approach 1. With OFF initially and the post-switch convention, the successive states are ON, OFF, ON, OFF,\ldots. They have no limit: arbitrarily close to midnight both states occur. There is no last switch before midnight. Time convergence does not imply state convergence.

The instructions for $t<1$ alone do not determine the state at $t=1$. Appending ON or OFF as a separate endpoint assignment is mathematically consistent with those earlier instructions. Requiring continuity at 1 is impossible. Do not offer ``half on'' as an endpoint deduction or treat infinity as an even/odd integer.

By contrast, a counter starting at 0 and moving halfway toward 1 has positions $1-2^{-n}$ converging to 1. Experiments can show initial marks; the recurring halfway rule explains arbitrarily small remaining gaps. Problem 3 turns both facts around. A card settles exactly when it flips at only finitely many marks, since each flip changes its side. A counter whose every move is at least some fixed length, for instance one that jumps end to end, cannot approach one place: near a place it approaches, all its later moves would be short. The converse fails, since short moves alone do not make a counter settle. Grades 2--3 manipulate, compare and design rules (Problems 1--3); Grades 4--5 can also separate an endpoint convention from a consequence (Problems 4--6). K--1 has an adult-led route on page 2. No actual supertask is performed.
\newpage
\section*{Preparation, readiness and first route}
Per pair: the time strip and the separate position strip cut from page 6, one counter at most 15 mm across, scissors, glue, pencil and ruler. Cut the two ON/OFF panels and glue them back-to-back to make one reversible card. Children fold the cut-out time strip and copy the folded marks onto the printed line in Problem 1. Work on the first four halves only, leaving a physically usable gap; describe or imagine later marks rather than folding them. Very late marks are mathematical descriptions, not demands for microscopic accurate folding. Put the two strips beside each other. Use the counter's center as its position; stop physical moves before counter size hides the gap. The endpoint remains visible and separate from switch marks.

Prerequisites: distinguish an endpoint from an interior mark; follow alternation; fold/locate halfway. Fraction notation is optional. One anchored adult reads the shared rule, witnesses each card flip and records a few predictions. The other child checks the move. No timer or flashing light is needed.

After the movement break, allow two minutes with strips/cards. In a three-minute whole-group launch, tell the 11:00--midnight light story, then demonstrate folding the gap from 11:30 to midnight to locate 11:45 and flipping ON to OFF. Then fold the START--END strip once and slide a counter from START to the fold. The worked visual uses these same clock labels. The separate START--END strip describes position, not time. Let children act before naming limits.

First route: for Grades 2--3 use Problems 1--3, with adult-scribed observations. Problems 4--6 are a Grades 4--5 continuation when endpoint discussion is ready. Allow 15 minutes for marks and flips; 10 minutes for halfway-counter motion on a fresh line; 10 minutes for comparing what settles and designing Problem 3's rules. A table that continues spends about 10 minutes on the endpoint rules. Preserve time for discussion. The rules are enforced by the partner checking marks and flipping the \emph{whole} two-state card. The optional endpoint conversation needs an adult; it is not a test of children saying the expected phrase. Timing and handling are untested.

K--1, adult-led on pages 1 and 6 (no K--1 pages are printed): the adult tells the light story and folds the cut-out time strip, and the child flips the card at each fold and says ON or OFF. After four folds, ask, ``Could we fold again? Will there ever be a last fold?'' Then fold the START--END strip and let the child slide the counter halfway toward END at each fold. Nothing needs reading. In plain words: there is always room to fold halfway again, so there is no last fold; the card keeps changing, so it never stays ON or stays OFF; the counter keeps getting closer to END but never lands on it. Stop when the folds get too small to make.
\newpage
\section*{Answers and conversation moves}
Problem 1: a farther mark always exists: between any switch mark and midnight put the next halfway mark. Thus the ``last drawn mark'' is not a last mark of the rule. The original switch-time sequence has a visible endpoint in its description but no terminal event.

Problem 3: a card settles if it flips at only finitely many marks, for example only at the first two (OFF, then ON, then OFF for ever after); never flipping counts too. A counter that jumps from one end to the other at every mark never approaches one place. A counter that stops moving stays at one place: like the card that never flips, it has settled, so it does not answer this question; ask for a rule that keeps the counter moving at every mark but never settles. Richer answers: flipping at every second mark never settles, because a card that flips at infinitely many marks keeps changing side; a counter that moves halfway toward END when the card turns ON and halfway toward START when it turns OFF goes 90, 45, 112.5, 56.25, \ldots\ mm and swings between about 60 and 120 mm for ever, every move at least 45 mm.

Problems 4--5: the switching rule fixes the card at every time before midnight and says nothing about midnight, so ON and OFF both fit; that is Problem 4's answer. The MIDNIGHT box holds the card or a written answer. Either state can be added as an extra rule without changing an earlier switch, and neither follows from the switching rule. If a child predicts ON at midnight, invite an OFF group to identify a different additional assumption. Ask whether either is forced by the earlier moves. If a child says undefined, ask what is specified before midnight and what information would finish the instruction sheet. Preserve the distinction between a missing endpoint rule and a nonexistent ordinary limit. A rule such as ``flip once more at midnight'' or ``stay as it was'' needs a last state before midnight, and there is none, so it does not choose a side; this is the same obstacle Problem 6 meets.

Problems 2 and 6: for the moving counter, after each action the distance remaining is half the old distance. Given any positive real tolerance, sufficiently many halvings make the gap smaller. The counter never arrives in a finite move, but its positions have endpoint 1 as their limit. This is not a proof from four drawings: the rule continues for every finite number of actions.

For the lamp, the proposed settled-earlier-state rule in Problem 6 fails. Take the odd and even post-switch subsequences; one is always ON and the other always OFF. They cannot converge to one common state. Averages of states answer a different question and do not turn the binary lamp half on.

Hints belong here: try predicting the next \emph{two} states; keep the state card separate from the timeline; compare remaining distance with the previous distance.
\newpage
\section*{Encore, assumptions and sources}
Encore: after Problem 3, design a counter rule whose positions approach the middle of the strip instead of END. Compare the number of actions, their times and their states separately. Return to the numbered-ball detective process: each label there does settle, even though the total count does not pass to the membership limit.

Assumptions: the two possible states are distinct; the rule covers all prescribed times strictly before midnight; a state immediately after a switch is used when describing the alternating sequence. Endpoint ON/OFF assignments do not claim physical realizability or continuity. No experiment with an actual infinitely fast mechanism is attempted.

Sources: James F. Thomson, ``Tasks and Super-Tasks'' (1954), especially the lamp discussion pp. 5--6; Paul Benacerraf, ``Tasks, Super-Tasks, and the Modern Eleatics'' (1962), pp. 768--770, distinguishing earlier behavior and endpoint specification. Original classroom presentation based on the organizer's example. Links and bibliographic details are in SOURCES.md.

Pedagogy: Rozhkovskaya, Lesson 7 ``At the lesson,'' item 1 notes missed legal moves during game play. Partner witnessing is our response, informed also by the organizer's Week 2 experience. This lamp activity uses folding and an actual reversible card.

Revised October 23, 2026 from the review card, after the organizer's October 17 approval: Problem 3 is new and the old Problems 3--5 are now 4--6. Unpiloted. US Letter, single-sided, Actual Size. Physical fit, folding tolerances, timing and classroom rehearsal untested. After use record actual marks attempted, children's endpoint proposals, whether they distinguished time/state, and where the adult supplied the investigation rather than read/recorded.
\end{document}
